\addcontentsline{toc}{section}{Einleitung}\section*{Einleitung}

Die Idee zu dieser Arbeit ist in Grundzügen schon früh entstanden. Im Januar 2020 während einer der ersten Vorlesungen der GymInf-Ausbildung, \textit{Einführung in Computersysteme} von Ulrich Ultes-Nitsche (Université de Fribourg) haben wir in einem Logikgatter-Simulator die Funktionsweise eines FlipFlops untersucht, einen adressierten Speicher modelliert, Halb- und Volladdierer entworfen und zu einem Addierer kombiniert. Diese Aufgaben haben mir viel Spass bereitet und einen ersten Einblick in den inneren Aufbau eines Computers ermöglicht. Zu diesem Zeitpunkt habe ich mich ein erstes Mal gefragt, ob es nicht möglich wäre, einen vollständigen Computer auf diese Art zu simulieren. In der Vorlesung haben wir das Thema jedoch nicht weiter vertieft und einen Sprung auf eine höhere Abstraktionsebene gemacht, in welcher die Funktionsweise von Betriebssystemen besprochen wurde. Nach meinem Empfinden blieb zwischen den NAND-Gattern und dem Linux-Kernel eine Lücke.

Die nächste Inspiration kam von einem Arbeitskollegen: Thomas Graf hat mich auf die Webseite von Ben Eater \cite{ben-eater} aufmerksam gemacht, der eine Anleitung liefert, wie man einen 8-Bit-Computer auf einem Steckbrett realisieren kann. Dabei folgt er einem \textit{simple-as-possible}-Aufbau, der von Paul Malvino in seinem Buch \enquote{Digital computer electronics} \cite{malvino} vorgeschlagen wurde. Das Projekt hat mich fasziniert. Würde ich damit die Lücke schliessen können? Erste zaghafte Versuche zeigten aber schnell, dass Ben Eaters Projekte zwar sehr spannend sind, aber viel Zeit und Energie darauf verwendet wird, immer wieder mit Problemen der physikalischen Welt wie schwankende Spannungen oder unsaubere Kontakte umzugehen. Der Fokus auf die logische Funktionsweise eines Computers geht damit ein Stück weit verloren. 

Den entscheidenden Hinweis, der schliesslich zu der vorliegenden Arbeit geführt hat, habe ich im Frühling 2021 von Matthias Hauswirth (Università della Svizzera italiana) erhalten. Er hat mich auf das nand2tetris-Projekt von Simon Schocken und Noam Nisan \cite{nand2tetris-url}, \cite{nand2tetris} aufmerksam gemacht, wofür ich ihm sehr dankbar bin! In den darauffolgenden Wochen habe ich voller Begeisterung die zwölf Projekte des nand2tetris-Kurses absolviert. Beginnend mit NAND-Gattern habe ich den Bau eines Computers in einer Hardwarebeschreibungssprache simuliert und anschliessend den Weg von einer einfachen Maschinensprache zu einer höheren Programmiersprache nachgezeichnet. Wie der Name des Kurses verrät, endet das Projekt damit, dass man das Spiel Tetris auf dem selbst gebauten Computer implementieren kann.

Damit war die Lücke geschlossen, und zwar auf eine enorm motivierende Art, die mit der grossen Befriedigung verbunden ist, einen eigenen Computer zu bauen - ganz ohne Lötkolben! Bereits während des Bearbeitens der einzelnen Projekte habe ich mir die Frage gestellt, ob es nicht möglich wäre, einen ähnlichen Kurs mit Gymnasialschülerinnen und Gymnasialschüler durchzuführen. Die zweite Hälfte des nand2tetris-Kurses, in welcher es um die Entwicklung einer höheren Programmiersprache auf der Basis der Maschinensprache geht, erfordert solide Programmierkenntnisse und liegt über dem Mittelschulniveau. Der erste Teil hingegen -- die Entwicklung des Computers aus NAND-Gattern -- schien mir, was das Thema und die erforderlichen Vorkenntnisse angeht, grundsätzlich für Gymnasialschülerinnen und Gymnasialschüler möglich. 

%Freilich würden einige Anpassungen nötig sein, denn das Werk von Schocken und Nisan richtet sich klar an Studierende. 

In der vorliegenden Arbeit wurden die Ideen von nand2tetris ergänzt, neu formuliert und so adaptiert, dass sie sich für den Einsatz im Rahmen des Ergänzungsfachs Informatik an einem Schweizer Gymnasium eignen. Anstelle der Hardwarebeschreibungssprache kommt mit CircuitVerse (\url{https://circuitverse.org}) ein Logikgatter-Simulator zum Einsatz. Dank dessen grafischer Oberfläche können die Schülerinnen und Schüler sprichwörtlich beobachten, wie ihre Schaltkreise funktionieren. Ausserdem wurden die Projekte, die zum Bau des Computers führen, durch Beispiele und klassische Aufgaben ergänzt, welche die Lernenden beim Aneignen der theoretischen Inhalte unterstützen. 

%Kurzum: Der in nand2tetris vorgestellte Weg von einfachen Logikgatter zu einem programmierbaren Computer bildet zwar immer noch den roten Faden dieser Unterrichtseinheit. Dieser Faden wurde aber ganz neu eingekleidet.